\documentclass[aspectratio=1610,11pt]{beamer}
% Shared style of the lecture slides (Slides/slidesNN.tex). Each deck starts with
%     \documentclass[aspectratio=1610,11pt]{beamer}
%     % Shared style of the lecture slides (Slides/slidesNN.tex). Each deck starts with
%     \documentclass[aspectratio=1610,11pt]{beamer}
%     % Shared style of the lecture slides (Slides/slidesNN.tex). Each deck starts with
%     \documentclass[aspectratio=1610,11pt]{beamer}
%     \input{slides_preamble}
% aspectratio=1610 (16:10) matches the Samsung Galaxy Tab S9 FE+ (2560 x 1600) used for teaching.
%
% The decks have no overlays (\pause, \only): every slide is one PDF page, so what you write on a
% page while sharing the tablet screen stays on that page. Board pages (\board) are white pages with
% a small header, for calculations and proofs written live.

\usepackage[T1]{fontenc}
\usepackage{lmodern}
\usefonttheme[onlymath]{serif}
\usepackage{amsmath,amssymb}
\usepackage{booktabs}
\usepackage{tikz}
\usetikzlibrary{arrows.meta,calc,positioning,shapes.geometric,decorations.pathreplacing}
\usepackage{pgfplots}
\pgfplotsset{compat=1.18}
\usepgfplotslibrary{fillbetween}

% Numbers of the lecture notes. A deck says which notes it follows with \setnotes{2} and
% loads their labels with \externaldocument[L2-]{../Lecture2/lecture2} (the .aux of the notes
% must exist). Then \begin{theorem}[Spectral theorem][thm:spectral] prints "notes L2, 10.1" at
% the right of the block title, and \secref{sec:spectral} prints "§2.10".
\usepackage{xr-hyper}
\makeatletter
\newcommand{\setnotes}[1]{\def\slides@notes{#1}}
\setnotes{2}
\newcommand{\notesref}[1]{\ref*{L\slides@notes-#1}}
\newcommand{\secref}[1]{\S\slides@notes.\ref*{L\slides@notes-#1}}
\makeatother

% Macros of the lecture notes.
\newcommand{\norm}[1]{\lVert #1 \rVert}
\newcommand{\ip}[2]{\langle #1, #2 \rangle}
\newcommand{\py}[1]{\texttt{\detokenize{#1}}}
\newcommand{\E}{\mathbb{E}}
\newcommand{\indep}{\perp\!\!\!\perp}
\newcommand{\R}{\mathbb{R}}
\newcommand{\bx}{\mathbf{x}}
\newcommand{\by}{\mathbf{y}}
\newcommand{\bw}{\mathbf{w}}
\newcommand{\bv}{\mathbf{v}}
\newcommand{\bu}{\mathbf{u}}
\newcommand{\bq}{\mathbf{q}}
\newcommand{\bb}{\mathbf{b}}
\newcommand{\be}{\mathbf{e}}
\DeclareMathOperator{\Var}{Var}
\DeclareMathOperator{\Cov}{Cov}
\DeclareMathOperator{\Corr}{Corr}
\DeclareMathOperator{\KL}{KL}
\DeclareMathOperator{\im}{im}
\DeclareMathOperator{\rank}{rank}
\DeclareMathOperator{\Null}{null}
\DeclareMathOperator{\col}{col}
\DeclareMathOperator{\row}{row}
\DeclareMathOperator{\Span}{span}
\DeclareMathOperator{\tr}{tr}
\DeclareMathOperator{\diag}{diag}
\DeclareMathOperator{\Div}{div}
\DeclareMathOperator{\curl}{curl}
\DeclareMathOperator{\vol}{vol}
\DeclareMathOperator{\area}{area}
\DeclareMathOperator{\sign}{sign}
\DeclareMathOperator{\softmax}{softmax}
\DeclareMathOperator{\ReLU}{ReLU}
\DeclareMathOperator*{\argmin}{arg\,min}
\DeclareMathOperator*{\argmax}{arg\,max}

% Colors.
\definecolor{accent}{RGB}{31,78,140}      % titles, theorems
\definecolor{defcolor}{RGB}{24,110,92}    % definitions
\definecolor{excolor}{RGB}{150,82,20}     % examples
\definecolor{mlcolor}{RGB}{120,40,120}    % machine learning connections
\definecolor{boardink}{RGB}{140,140,140}  % header of the board pages

% Layout.
\setbeamertemplate{navigation symbols}{}
\setbeamersize{text margin left=0.9em, text margin right=0.9em}
\setbeamercolor{structure}{fg=accent}
\setbeamercolor{normal text}{fg=black, bg=white}
\setbeamercolor{frametitle}{fg=accent, bg=white}
\setbeamerfont{frametitle}{series=\bfseries, size=\large}
\setbeamertemplate{frametitle}{%
    \vskip0.45em\insertframetitle\par\vskip-0.25em
    {\color{accent!25}\rule{\linewidth}{0.5pt}}\vskip-0.2em}
\setbeamertemplate{itemize item}{\color{accent}\raisebox{0.15ex}{\scriptsize$\blacktriangleright$}}
\setbeamertemplate{itemize subitem}{\color{accent}\raisebox{0.2ex}{\tiny$\bullet$}}
\setbeamertemplate{enumerate item}{\color{accent}\insertenumlabel.}
\setbeamerfont{footline}{size=\tiny}
\setbeamertemplate{footline}{%
    \leavevmode\hbox to \paperwidth{\hskip0.9em{\color{gray}\insertshorttitle}\hfill
    {\color{gray}\insertframenumber}\hskip0.9em}\vskip0.45em}
\setbeamertemplate{blocks}[rounded][shadow=false]
\setbeamercolor{block title}{fg=white, bg=accent}
\setbeamercolor{block body}{fg=black, bg=accent!6}
\setbeamerfont{block title}{size=\normalsize, series=\bfseries}
\setlength{\abovedisplayskip}{4pt}
\setlength{\belowdisplayskip}{4pt}

% Blocks: \begin{theorem}[Spectral theorem][thm:spectral] ... \end{theorem}; both arguments are optional.
\makeatletter
\newcommand{\slides@block}[5]{% color, word, name, label in the notes, body tint
    \setbeamercolor{block title}{fg=white, bg=#1}%
    \setbeamercolor{block body}{fg=black, bg=#1!#5}%
    \begin{block}{#2\if\relax\detokenize{#3}\relax\else\ {\normalfont(#3)}\fi
        \if\relax\detokenize{#4}\relax\else\hfill{\normalfont\scriptsize notes~L\slides@notes, \notesref{#4}}\fi}}
\DeclareDocumentEnvironment{theorem}{O{} O{}}{\slides@block{accent}{Theorem}{#1}{#2}{6}}{\end{block}}
\DeclareDocumentEnvironment{proposition}{O{} O{}}{\slides@block{accent}{Proposition}{#1}{#2}{6}}{\end{block}}
\DeclareDocumentEnvironment{lemma}{O{} O{}}{\slides@block{accent}{Lemma}{#1}{#2}{6}}{\end{block}}
\DeclareDocumentEnvironment{corollary}{O{} O{}}{\slides@block{accent}{Corollary}{#1}{#2}{6}}{\end{block}}
\DeclareDocumentEnvironment{definition}{O{} O{}}{\slides@block{defcolor}{Definition}{#1}{#2}{6}}{\end{block}}
\DeclareDocumentEnvironment{example}{O{} O{}}{\slides@block{excolor}{Example}{#1}{#2}{6}}{\end{block}}
\DeclareDocumentEnvironment{remark}{O{} O{}}{\slides@block{gray!70!black}{Remark}{#1}{#2}{5}}{\end{block}}
\DeclareDocumentEnvironment{mlbox}{O{In machine learning}}{%
    \setbeamercolor{block title}{fg=white, bg=mlcolor}%
    \setbeamercolor{block body}{fg=black, bg=mlcolor!6}%
    \begin{block}{#1}}{\end{block}}
\makeatother

% Board pages.
%   \board{Title}{Task, shown in small print at the top}            one white page
%   \board[3]{Title}{Task}                                          three pages (the last two blank)
\newcounter{boardpage}
\newcommand{\boardhead}[1]{%
    \vskip0.3em
    \noindent{\color{boardink}\scriptsize
        \tikz[baseline=(b.base)]\node[draw=boardink, rounded corners=2pt, inner sep=1.5pt,
            font=\tiny\bfseries\sffamily] (b) {BOARD};\ \ #1}\par\vskip0.3em}
\newcommand{\board}[3][1]{%
    \begin{frame}[t]
        \boardhead{#2}
        \if\relax\detokenize{#3}\relax\else
            {\small #3\par}
        \fi
    \end{frame}%
    \setcounter{boardpage}{1}%
    \loop\ifnum\value{boardpage}<#1\relax
        \stepcounter{boardpage}%
        \begin{frame}[t]\boardhead{#2 (continued)}\end{frame}%
    \repeat}

% Practice pages: problems of the practice file, as white pages for solving them live.
% A deck loads the practice file's labels with \externaldocument[P1-]{../Lecture1/practice1} and
% \setpractice{1}. Then
%   \practice{pr:eig2}{\mandatory}{Problem text}    one page, headed "PRACTICE  Problem 10.1 (mandatory)"
%   \practice[2]{pr:eig-props}{}{Problem text}      two pages (the second blank)
%   \practice{pr:eig-gre}{$\star$}{Problem text}    the marks of the practice file
% The answers stay in LectureN/practiceN_answers.tex.
\newcommand{\mandatory}{\textcolor{red!70!black}{\textbf{mandatory}}}
\makeatletter
\newcommand{\setpractice}[1]{\def\slides@practice{#1}}
\newcommand{\practiceref}[1]{\ref*{P\slides@practice-#1}}
\newcommand{\practicehead}[2]{% label in the practice file, text after the problem number
    \vskip0.3em
    \noindent{\color{boardink}\scriptsize
        \tikz[baseline=(b.base)]\node[draw=defcolor, text=defcolor, rounded corners=2pt, inner sep=1.5pt,
            font=\tiny\bfseries\sffamily] (b) {PRACTICE};\ \ Practice~\slides@practice, Problem~\practiceref{#1}#2}%
    \par\vskip0.3em}
\newcommand{\practice}[4][1]{%
    \begin{frame}[t]
        \practicehead{#2}{\if\relax\detokenize{#3}\relax\else\ (#3)\fi}
        {\small #4\par}
    \end{frame}%
    \setcounter{boardpage}{1}%
    \loop\ifnum\value{boardpage}<#1\relax
        \stepcounter{boardpage}%
        \begin{frame}[t]\practicehead{#2}{\ (continued)}\end{frame}%
    \repeat}
\makeatother

% Divider page: \topic{\secref{sec:spectral}}{The spectral theorem} or \topic{\S2.1--2.2}{Norms}.
\newcommand{\topic}[2]{%
    \begin{frame}[c]
        \centering
        {\color{gray}\large #1}\par\vskip0.4em
        {\color{accent}\LARGE\bfseries #2}\par
    \end{frame}}

% Title page.
\setbeamercolor{title}{fg=accent}
\setbeamerfont{title}{series=\bfseries, size=\Large}
\setbeamercolor{subtitle}{fg=gray!60!black}
\setbeamerfont{subtitle}{size=\normalsize}
\setbeamertemplate{title page}{%
    \vfill
    {\usebeamercolor[fg]{title}\usebeamerfont{title}\inserttitle\par}
    \vskip0.6em
    {\usebeamercolor[fg]{subtitle}\usebeamerfont{subtitle}\insertsubtitle\par}
    \vskip1.2em
    {\small\color{gray}\insertauthor\par}
    \vfill}
\author{Mathematics for Machine Learning}
\date{}

% Two columns with a picture: \sidebyside{0.5}{left text}{right picture}.
\newcommand{\sidebyside}[3]{%
    \begin{columns}[T]
        \begin{column}{#1\linewidth}#2\end{column}
        \begin{column}{\dimexpr\linewidth-#1\linewidth-1em\relax}#3\end{column}
    \end{columns}}

% "Check" slide after a board computation.
\newcommand{\checkslide}[2]{%
    \begin{frame}{Check: #1}
        #2
    \end{frame}}

% Pictures.
\newcommand{\notesfig}[2][1]{\resizebox{#1\linewidth}{!}{\input{figures/notes/#2}}}
\newcommand{\fig}[2][1]{\includegraphics[width=#1\linewidth]{figures/#2}}

% aspectratio=1610 (16:10) matches the Samsung Galaxy Tab S9 FE+ (2560 x 1600) used for teaching.
%
% The decks have no overlays (\pause, \only): every slide is one PDF page, so what you write on a
% page while sharing the tablet screen stays on that page. Board pages (\board) are white pages with
% a small header, for calculations and proofs written live.

\usepackage[T1]{fontenc}
\usepackage{lmodern}
\usefonttheme[onlymath]{serif}
\usepackage{amsmath,amssymb}
\usepackage{booktabs}
\usepackage{tikz}
\usetikzlibrary{arrows.meta,calc,positioning,shapes.geometric,decorations.pathreplacing}
\usepackage{pgfplots}
\pgfplotsset{compat=1.18}
\usepgfplotslibrary{fillbetween}

% Numbers of the lecture notes. A deck says which notes it follows with \setnotes{2} and
% loads their labels with \externaldocument[L2-]{../Lecture2/lecture2} (the .aux of the notes
% must exist). Then \begin{theorem}[Spectral theorem][thm:spectral] prints "notes L2, 10.1" at
% the right of the block title, and \secref{sec:spectral} prints "§2.10".
\usepackage{xr-hyper}
\makeatletter
\newcommand{\setnotes}[1]{\def\slides@notes{#1}}
\setnotes{2}
\newcommand{\notesref}[1]{\ref*{L\slides@notes-#1}}
\newcommand{\secref}[1]{\S\slides@notes.\ref*{L\slides@notes-#1}}
\makeatother

% Macros of the lecture notes.
\newcommand{\norm}[1]{\lVert #1 \rVert}
\newcommand{\ip}[2]{\langle #1, #2 \rangle}
\newcommand{\py}[1]{\texttt{\detokenize{#1}}}
\newcommand{\E}{\mathbb{E}}
\newcommand{\indep}{\perp\!\!\!\perp}
\newcommand{\R}{\mathbb{R}}
\newcommand{\bx}{\mathbf{x}}
\newcommand{\by}{\mathbf{y}}
\newcommand{\bw}{\mathbf{w}}
\newcommand{\bv}{\mathbf{v}}
\newcommand{\bu}{\mathbf{u}}
\newcommand{\bq}{\mathbf{q}}
\newcommand{\bb}{\mathbf{b}}
\newcommand{\be}{\mathbf{e}}
\DeclareMathOperator{\Var}{Var}
\DeclareMathOperator{\Cov}{Cov}
\DeclareMathOperator{\Corr}{Corr}
\DeclareMathOperator{\KL}{KL}
\DeclareMathOperator{\im}{im}
\DeclareMathOperator{\rank}{rank}
\DeclareMathOperator{\Null}{null}
\DeclareMathOperator{\col}{col}
\DeclareMathOperator{\row}{row}
\DeclareMathOperator{\Span}{span}
\DeclareMathOperator{\tr}{tr}
\DeclareMathOperator{\diag}{diag}
\DeclareMathOperator{\Div}{div}
\DeclareMathOperator{\curl}{curl}
\DeclareMathOperator{\vol}{vol}
\DeclareMathOperator{\area}{area}
\DeclareMathOperator{\sign}{sign}
\DeclareMathOperator{\softmax}{softmax}
\DeclareMathOperator{\ReLU}{ReLU}
\DeclareMathOperator*{\argmin}{arg\,min}
\DeclareMathOperator*{\argmax}{arg\,max}

% Colors.
\definecolor{accent}{RGB}{31,78,140}      % titles, theorems
\definecolor{defcolor}{RGB}{24,110,92}    % definitions
\definecolor{excolor}{RGB}{150,82,20}     % examples
\definecolor{mlcolor}{RGB}{120,40,120}    % machine learning connections
\definecolor{boardink}{RGB}{140,140,140}  % header of the board pages

% Layout.
\setbeamertemplate{navigation symbols}{}
\setbeamersize{text margin left=0.9em, text margin right=0.9em}
\setbeamercolor{structure}{fg=accent}
\setbeamercolor{normal text}{fg=black, bg=white}
\setbeamercolor{frametitle}{fg=accent, bg=white}
\setbeamerfont{frametitle}{series=\bfseries, size=\large}
\setbeamertemplate{frametitle}{%
    \vskip0.45em\insertframetitle\par\vskip-0.25em
    {\color{accent!25}\rule{\linewidth}{0.5pt}}\vskip-0.2em}
\setbeamertemplate{itemize item}{\color{accent}\raisebox{0.15ex}{\scriptsize$\blacktriangleright$}}
\setbeamertemplate{itemize subitem}{\color{accent}\raisebox{0.2ex}{\tiny$\bullet$}}
\setbeamertemplate{enumerate item}{\color{accent}\insertenumlabel.}
\setbeamerfont{footline}{size=\tiny}
\setbeamertemplate{footline}{%
    \leavevmode\hbox to \paperwidth{\hskip0.9em{\color{gray}\insertshorttitle}\hfill
    {\color{gray}\insertframenumber}\hskip0.9em}\vskip0.45em}
\setbeamertemplate{blocks}[rounded][shadow=false]
\setbeamercolor{block title}{fg=white, bg=accent}
\setbeamercolor{block body}{fg=black, bg=accent!6}
\setbeamerfont{block title}{size=\normalsize, series=\bfseries}
\setlength{\abovedisplayskip}{4pt}
\setlength{\belowdisplayskip}{4pt}

% Blocks: \begin{theorem}[Spectral theorem][thm:spectral] ... \end{theorem}; both arguments are optional.
\makeatletter
\newcommand{\slides@block}[5]{% color, word, name, label in the notes, body tint
    \setbeamercolor{block title}{fg=white, bg=#1}%
    \setbeamercolor{block body}{fg=black, bg=#1!#5}%
    \begin{block}{#2\if\relax\detokenize{#3}\relax\else\ {\normalfont(#3)}\fi
        \if\relax\detokenize{#4}\relax\else\hfill{\normalfont\scriptsize notes~L\slides@notes, \notesref{#4}}\fi}}
\DeclareDocumentEnvironment{theorem}{O{} O{}}{\slides@block{accent}{Theorem}{#1}{#2}{6}}{\end{block}}
\DeclareDocumentEnvironment{proposition}{O{} O{}}{\slides@block{accent}{Proposition}{#1}{#2}{6}}{\end{block}}
\DeclareDocumentEnvironment{lemma}{O{} O{}}{\slides@block{accent}{Lemma}{#1}{#2}{6}}{\end{block}}
\DeclareDocumentEnvironment{corollary}{O{} O{}}{\slides@block{accent}{Corollary}{#1}{#2}{6}}{\end{block}}
\DeclareDocumentEnvironment{definition}{O{} O{}}{\slides@block{defcolor}{Definition}{#1}{#2}{6}}{\end{block}}
\DeclareDocumentEnvironment{example}{O{} O{}}{\slides@block{excolor}{Example}{#1}{#2}{6}}{\end{block}}
\DeclareDocumentEnvironment{remark}{O{} O{}}{\slides@block{gray!70!black}{Remark}{#1}{#2}{5}}{\end{block}}
\DeclareDocumentEnvironment{mlbox}{O{In machine learning}}{%
    \setbeamercolor{block title}{fg=white, bg=mlcolor}%
    \setbeamercolor{block body}{fg=black, bg=mlcolor!6}%
    \begin{block}{#1}}{\end{block}}
\makeatother

% Board pages.
%   \board{Title}{Task, shown in small print at the top}            one white page
%   \board[3]{Title}{Task}                                          three pages (the last two blank)
\newcounter{boardpage}
\newcommand{\boardhead}[1]{%
    \vskip0.3em
    \noindent{\color{boardink}\scriptsize
        \tikz[baseline=(b.base)]\node[draw=boardink, rounded corners=2pt, inner sep=1.5pt,
            font=\tiny\bfseries\sffamily] (b) {BOARD};\ \ #1}\par\vskip0.3em}
\newcommand{\board}[3][1]{%
    \begin{frame}[t]
        \boardhead{#2}
        \if\relax\detokenize{#3}\relax\else
            {\small #3\par}
        \fi
    \end{frame}%
    \setcounter{boardpage}{1}%
    \loop\ifnum\value{boardpage}<#1\relax
        \stepcounter{boardpage}%
        \begin{frame}[t]\boardhead{#2 (continued)}\end{frame}%
    \repeat}

% Practice pages: problems of the practice file, as white pages for solving them live.
% A deck loads the practice file's labels with \externaldocument[P1-]{../Lecture1/practice1} and
% \setpractice{1}. Then
%   \practice{pr:eig2}{\mandatory}{Problem text}    one page, headed "PRACTICE  Problem 10.1 (mandatory)"
%   \practice[2]{pr:eig-props}{}{Problem text}      two pages (the second blank)
%   \practice{pr:eig-gre}{$\star$}{Problem text}    the marks of the practice file
% The answers stay in LectureN/practiceN_answers.tex.
\newcommand{\mandatory}{\textcolor{red!70!black}{\textbf{mandatory}}}
\makeatletter
\newcommand{\setpractice}[1]{\def\slides@practice{#1}}
\newcommand{\practiceref}[1]{\ref*{P\slides@practice-#1}}
\newcommand{\practicehead}[2]{% label in the practice file, text after the problem number
    \vskip0.3em
    \noindent{\color{boardink}\scriptsize
        \tikz[baseline=(b.base)]\node[draw=defcolor, text=defcolor, rounded corners=2pt, inner sep=1.5pt,
            font=\tiny\bfseries\sffamily] (b) {PRACTICE};\ \ Practice~\slides@practice, Problem~\practiceref{#1}#2}%
    \par\vskip0.3em}
\newcommand{\practice}[4][1]{%
    \begin{frame}[t]
        \practicehead{#2}{\if\relax\detokenize{#3}\relax\else\ (#3)\fi}
        {\small #4\par}
    \end{frame}%
    \setcounter{boardpage}{1}%
    \loop\ifnum\value{boardpage}<#1\relax
        \stepcounter{boardpage}%
        \begin{frame}[t]\practicehead{#2}{\ (continued)}\end{frame}%
    \repeat}
\makeatother

% Divider page: \topic{\secref{sec:spectral}}{The spectral theorem} or \topic{\S2.1--2.2}{Norms}.
\newcommand{\topic}[2]{%
    \begin{frame}[c]
        \centering
        {\color{gray}\large #1}\par\vskip0.4em
        {\color{accent}\LARGE\bfseries #2}\par
    \end{frame}}

% Title page.
\setbeamercolor{title}{fg=accent}
\setbeamerfont{title}{series=\bfseries, size=\Large}
\setbeamercolor{subtitle}{fg=gray!60!black}
\setbeamerfont{subtitle}{size=\normalsize}
\setbeamertemplate{title page}{%
    \vfill
    {\usebeamercolor[fg]{title}\usebeamerfont{title}\inserttitle\par}
    \vskip0.6em
    {\usebeamercolor[fg]{subtitle}\usebeamerfont{subtitle}\insertsubtitle\par}
    \vskip1.2em
    {\small\color{gray}\insertauthor\par}
    \vfill}
\author{Mathematics for Machine Learning}
\date{}

% Two columns with a picture: \sidebyside{0.5}{left text}{right picture}.
\newcommand{\sidebyside}[3]{%
    \begin{columns}[T]
        \begin{column}{#1\linewidth}#2\end{column}
        \begin{column}{\dimexpr\linewidth-#1\linewidth-1em\relax}#3\end{column}
    \end{columns}}

% "Check" slide after a board computation.
\newcommand{\checkslide}[2]{%
    \begin{frame}{Check: #1}
        #2
    \end{frame}}

% Pictures.
\newcommand{\notesfig}[2][1]{\resizebox{#1\linewidth}{!}{\input{figures/notes/#2}}}
\newcommand{\fig}[2][1]{\includegraphics[width=#1\linewidth]{figures/#2}}

% aspectratio=1610 (16:10) matches the Samsung Galaxy Tab S9 FE+ (2560 x 1600) used for teaching.
%
% The decks have no overlays (\pause, \only): every slide is one PDF page, so what you write on a
% page while sharing the tablet screen stays on that page. Board pages (\board) are white pages with
% a small header, for calculations and proofs written live.

\usepackage[T1]{fontenc}
\usepackage{lmodern}
\usefonttheme[onlymath]{serif}
\usepackage{amsmath,amssymb}
\usepackage{booktabs}
\usepackage{tikz}
\usetikzlibrary{arrows.meta,calc,positioning,shapes.geometric,decorations.pathreplacing}
\usepackage{pgfplots}
\pgfplotsset{compat=1.18}
\usepgfplotslibrary{fillbetween}

% Numbers of the lecture notes. A deck says which notes it follows with \setnotes{2} and
% loads their labels with \externaldocument[L2-]{../Lecture2/lecture2} (the .aux of the notes
% must exist). Then \begin{theorem}[Spectral theorem][thm:spectral] prints "notes L2, 10.1" at
% the right of the block title, and \secref{sec:spectral} prints "§2.10".
\usepackage{xr-hyper}
\makeatletter
\newcommand{\setnotes}[1]{\def\slides@notes{#1}}
\setnotes{2}
\newcommand{\notesref}[1]{\ref*{L\slides@notes-#1}}
\newcommand{\secref}[1]{\S\slides@notes.\ref*{L\slides@notes-#1}}
\makeatother

% Macros of the lecture notes.
\newcommand{\norm}[1]{\lVert #1 \rVert}
\newcommand{\ip}[2]{\langle #1, #2 \rangle}
\newcommand{\py}[1]{\texttt{\detokenize{#1}}}
\newcommand{\E}{\mathbb{E}}
\newcommand{\indep}{\perp\!\!\!\perp}
\newcommand{\R}{\mathbb{R}}
\newcommand{\bx}{\mathbf{x}}
\newcommand{\by}{\mathbf{y}}
\newcommand{\bw}{\mathbf{w}}
\newcommand{\bv}{\mathbf{v}}
\newcommand{\bu}{\mathbf{u}}
\newcommand{\bq}{\mathbf{q}}
\newcommand{\bb}{\mathbf{b}}
\newcommand{\be}{\mathbf{e}}
\DeclareMathOperator{\Var}{Var}
\DeclareMathOperator{\Cov}{Cov}
\DeclareMathOperator{\Corr}{Corr}
\DeclareMathOperator{\KL}{KL}
\DeclareMathOperator{\im}{im}
\DeclareMathOperator{\rank}{rank}
\DeclareMathOperator{\Null}{null}
\DeclareMathOperator{\col}{col}
\DeclareMathOperator{\row}{row}
\DeclareMathOperator{\Span}{span}
\DeclareMathOperator{\tr}{tr}
\DeclareMathOperator{\diag}{diag}
\DeclareMathOperator{\Div}{div}
\DeclareMathOperator{\curl}{curl}
\DeclareMathOperator{\vol}{vol}
\DeclareMathOperator{\area}{area}
\DeclareMathOperator{\sign}{sign}
\DeclareMathOperator{\softmax}{softmax}
\DeclareMathOperator{\ReLU}{ReLU}
\DeclareMathOperator*{\argmin}{arg\,min}
\DeclareMathOperator*{\argmax}{arg\,max}

% Colors.
\definecolor{accent}{RGB}{31,78,140}      % titles, theorems
\definecolor{defcolor}{RGB}{24,110,92}    % definitions
\definecolor{excolor}{RGB}{150,82,20}     % examples
\definecolor{mlcolor}{RGB}{120,40,120}    % machine learning connections
\definecolor{boardink}{RGB}{140,140,140}  % header of the board pages

% Layout.
\setbeamertemplate{navigation symbols}{}
\setbeamersize{text margin left=0.9em, text margin right=0.9em}
\setbeamercolor{structure}{fg=accent}
\setbeamercolor{normal text}{fg=black, bg=white}
\setbeamercolor{frametitle}{fg=accent, bg=white}
\setbeamerfont{frametitle}{series=\bfseries, size=\large}
\setbeamertemplate{frametitle}{%
    \vskip0.45em\insertframetitle\par\vskip-0.25em
    {\color{accent!25}\rule{\linewidth}{0.5pt}}\vskip-0.2em}
\setbeamertemplate{itemize item}{\color{accent}\raisebox{0.15ex}{\scriptsize$\blacktriangleright$}}
\setbeamertemplate{itemize subitem}{\color{accent}\raisebox{0.2ex}{\tiny$\bullet$}}
\setbeamertemplate{enumerate item}{\color{accent}\insertenumlabel.}
\setbeamerfont{footline}{size=\tiny}
\setbeamertemplate{footline}{%
    \leavevmode\hbox to \paperwidth{\hskip0.9em{\color{gray}\insertshorttitle}\hfill
    {\color{gray}\insertframenumber}\hskip0.9em}\vskip0.45em}
\setbeamertemplate{blocks}[rounded][shadow=false]
\setbeamercolor{block title}{fg=white, bg=accent}
\setbeamercolor{block body}{fg=black, bg=accent!6}
\setbeamerfont{block title}{size=\normalsize, series=\bfseries}
\setlength{\abovedisplayskip}{4pt}
\setlength{\belowdisplayskip}{4pt}

% Blocks: \begin{theorem}[Spectral theorem][thm:spectral] ... \end{theorem}; both arguments are optional.
\makeatletter
\newcommand{\slides@block}[5]{% color, word, name, label in the notes, body tint
    \setbeamercolor{block title}{fg=white, bg=#1}%
    \setbeamercolor{block body}{fg=black, bg=#1!#5}%
    \begin{block}{#2\if\relax\detokenize{#3}\relax\else\ {\normalfont(#3)}\fi
        \if\relax\detokenize{#4}\relax\else\hfill{\normalfont\scriptsize notes~L\slides@notes, \notesref{#4}}\fi}}
\DeclareDocumentEnvironment{theorem}{O{} O{}}{\slides@block{accent}{Theorem}{#1}{#2}{6}}{\end{block}}
\DeclareDocumentEnvironment{proposition}{O{} O{}}{\slides@block{accent}{Proposition}{#1}{#2}{6}}{\end{block}}
\DeclareDocumentEnvironment{lemma}{O{} O{}}{\slides@block{accent}{Lemma}{#1}{#2}{6}}{\end{block}}
\DeclareDocumentEnvironment{corollary}{O{} O{}}{\slides@block{accent}{Corollary}{#1}{#2}{6}}{\end{block}}
\DeclareDocumentEnvironment{definition}{O{} O{}}{\slides@block{defcolor}{Definition}{#1}{#2}{6}}{\end{block}}
\DeclareDocumentEnvironment{example}{O{} O{}}{\slides@block{excolor}{Example}{#1}{#2}{6}}{\end{block}}
\DeclareDocumentEnvironment{remark}{O{} O{}}{\slides@block{gray!70!black}{Remark}{#1}{#2}{5}}{\end{block}}
\DeclareDocumentEnvironment{mlbox}{O{In machine learning}}{%
    \setbeamercolor{block title}{fg=white, bg=mlcolor}%
    \setbeamercolor{block body}{fg=black, bg=mlcolor!6}%
    \begin{block}{#1}}{\end{block}}
\makeatother

% Board pages.
%   \board{Title}{Task, shown in small print at the top}            one white page
%   \board[3]{Title}{Task}                                          three pages (the last two blank)
\newcounter{boardpage}
\newcommand{\boardhead}[1]{%
    \vskip0.3em
    \noindent{\color{boardink}\scriptsize
        \tikz[baseline=(b.base)]\node[draw=boardink, rounded corners=2pt, inner sep=1.5pt,
            font=\tiny\bfseries\sffamily] (b) {BOARD};\ \ #1}\par\vskip0.3em}
\newcommand{\board}[3][1]{%
    \begin{frame}[t]
        \boardhead{#2}
        \if\relax\detokenize{#3}\relax\else
            {\small #3\par}
        \fi
    \end{frame}%
    \setcounter{boardpage}{1}%
    \loop\ifnum\value{boardpage}<#1\relax
        \stepcounter{boardpage}%
        \begin{frame}[t]\boardhead{#2 (continued)}\end{frame}%
    \repeat}

% Practice pages: problems of the practice file, as white pages for solving them live.
% A deck loads the practice file's labels with \externaldocument[P1-]{../Lecture1/practice1} and
% \setpractice{1}. Then
%   \practice{pr:eig2}{\mandatory}{Problem text}    one page, headed "PRACTICE  Problem 10.1 (mandatory)"
%   \practice[2]{pr:eig-props}{}{Problem text}      two pages (the second blank)
%   \practice{pr:eig-gre}{$\star$}{Problem text}    the marks of the practice file
% The answers stay in LectureN/practiceN_answers.tex.
\newcommand{\mandatory}{\textcolor{red!70!black}{\textbf{mandatory}}}
\makeatletter
\newcommand{\setpractice}[1]{\def\slides@practice{#1}}
\newcommand{\practiceref}[1]{\ref*{P\slides@practice-#1}}
\newcommand{\practicehead}[2]{% label in the practice file, text after the problem number
    \vskip0.3em
    \noindent{\color{boardink}\scriptsize
        \tikz[baseline=(b.base)]\node[draw=defcolor, text=defcolor, rounded corners=2pt, inner sep=1.5pt,
            font=\tiny\bfseries\sffamily] (b) {PRACTICE};\ \ Practice~\slides@practice, Problem~\practiceref{#1}#2}%
    \par\vskip0.3em}
\newcommand{\practice}[4][1]{%
    \begin{frame}[t]
        \practicehead{#2}{\if\relax\detokenize{#3}\relax\else\ (#3)\fi}
        {\small #4\par}
    \end{frame}%
    \setcounter{boardpage}{1}%
    \loop\ifnum\value{boardpage}<#1\relax
        \stepcounter{boardpage}%
        \begin{frame}[t]\practicehead{#2}{\ (continued)}\end{frame}%
    \repeat}
\makeatother

% Divider page: \topic{\secref{sec:spectral}}{The spectral theorem} or \topic{\S2.1--2.2}{Norms}.
\newcommand{\topic}[2]{%
    \begin{frame}[c]
        \centering
        {\color{gray}\large #1}\par\vskip0.4em
        {\color{accent}\LARGE\bfseries #2}\par
    \end{frame}}

% Title page.
\setbeamercolor{title}{fg=accent}
\setbeamerfont{title}{series=\bfseries, size=\Large}
\setbeamercolor{subtitle}{fg=gray!60!black}
\setbeamerfont{subtitle}{size=\normalsize}
\setbeamertemplate{title page}{%
    \vfill
    {\usebeamercolor[fg]{title}\usebeamerfont{title}\inserttitle\par}
    \vskip0.6em
    {\usebeamercolor[fg]{subtitle}\usebeamerfont{subtitle}\insertsubtitle\par}
    \vskip1.2em
    {\small\color{gray}\insertauthor\par}
    \vfill}
\author{Mathematics for Machine Learning}
\date{}

% Two columns with a picture: \sidebyside{0.5}{left text}{right picture}.
\newcommand{\sidebyside}[3]{%
    \begin{columns}[T]
        \begin{column}{#1\linewidth}#2\end{column}
        \begin{column}{\dimexpr\linewidth-#1\linewidth-1em\relax}#3\end{column}
    \end{columns}}

% "Check" slide after a board computation.
\newcommand{\checkslide}[2]{%
    \begin{frame}{Check: #1}
        #2
    \end{frame}}

% Pictures.
\newcommand{\notesfig}[2][1]{\resizebox{#1\linewidth}{!}{\input{figures/notes/#2}}}
\newcommand{\fig}[2][1]{\includegraphics[width=#1\linewidth]{figures/#2}}

\externaldocument[L3-]{../Lecture3/lecture3}
\setnotes{3}

\title[Lecture 7 · Calculus I]{Lecture 7. Calculus I:\\ Derivatives, the Chain Rule and Backpropagation}
\subtitle{Notes: Lecture 3, \S3.1--3.8 and \S3.15.2\\ Self-study: \S3.13--3.14, \S3.15.1, \S3.16}

\begin{document}

\begin{frame}[plain]
    \titlepage
\end{frame}

\begin{frame}{Today}
    \begin{columns}[T]
        \begin{column}{0.52\linewidth}
            \begin{enumerate}\small
                \item Functions of several variables (review of one variable, Taylor)
                \item Partial derivatives, gradients, differentiability
                \item Jacobians and the chain rule
                \item Gradients with respect to matrices
                \item Automatic differentiation and backpropagation
                \item Hessians and Taylor's theorem
                \item Gradient flow; gradient descent $=$ Euler's method
            \end{enumerate}
        \end{column}
        \begin{column}{0.44\linewidth}
            \begin{mlbox}[Used in]
                \small
                \begin{itemize}
                    \item the derivation of backpropagation (DL)
                    \item autograd in PyTorch (DL)
                    \item vanishing and exploding gradients (DL)
                \end{itemize}
            \end{mlbox}
        \end{column}
    \end{columns}
\end{frame}

\begin{frame}{Why derivatives?}
    \small
    A model has parameters $\boldsymbol{\theta}$ (millions of weights). A loss $L(\boldsymbol{\theta})$ measures how badly it fits the data. Training $=$ making $L$ small, step by step:
    \[\boldsymbol{\theta}_{k+1} = \boldsymbol{\theta}_k - \eta\,\nabla L(\boldsymbol{\theta}_k).\]
    We need:
    \begin{itemize}
        \item derivatives of functions $\R^n \to \R^m$: gradients and Jacobians;
        \item the chain rule: a network is a composition of layers;
        \item gradients with respect to weight \emph{matrices};
        \item an algorithm that computes all of them at the cost of a few evaluations of $L$: \textbf{backpropagation};
        \item second derivatives to understand the shape of $L$ near a point (and next time: optimization).
    \end{itemize}
    Linear algebra everywhere: derivatives are linear maps (matrices), steepest ascent is Cauchy--Schwarz, and the shape near a minimum is given by the eigenvalues of a symmetric matrix.
\end{frame}

% ------------------------------------------------------------------ Functions
\topic{\S3.1--3.2}{Functions of several variables}

\begin{frame}{Functions in machine learning}
    \small
    $f: \R^n \to \R^m$, \ $f(\bx) = (f_1(\bx), \ldots, f_m(\bx))^\top$.
    \begin{example}[][ex:ml-functions]
        \small
        \begin{itemize}
            \item linear model $f(\bx) = \bw^\top\bx + b$;
            \item a layer $f(\bx) = \sigma(W\bx + \mathbf{b})$, $\sigma$ applied entrywise: sigmoid $\frac{1}{1 + e^{-t}}$, $\tanh$, $\ReLU(t) = \max(0, t)$;
            \item softmax $\softmax(\mathbf{z})_i = \frac{e^{z_i}}{e^{z_1} + \cdots + e^{z_K}}$: scores $\to$ probabilities;
            \item a loss $L(\boldsymbol{\theta}) = \frac{1}{N}\sum_i \ell(f_{\boldsymbol{\theta}}(\bx_i), y_i)$, a function of the parameters.
        \end{itemize}
    \end{example}
    \centering\notesfig[0.64]{L3-level-sets}
\end{frame}

\begin{frame}{Limits and continuity}
    \begin{definition}[][def:limit]
        \small $\lim_{\bx \to \mathbf{a}} f(\bx) = \mathbf{L}$ if for every $\varepsilon > 0$ there is $\delta > 0$ with $\norm{f(\bx) - \mathbf{L}} < \varepsilon$ whenever $0 < \norm{\bx - \mathbf{a}} < \delta$. \ $f$ is continuous at $\mathbf{a}$ if $\lim_{\bx \to \mathbf{a}} f(\bx) = f(\mathbf{a})$.
    \end{definition}
    \begin{example}[][ex:xy]
        \small $f(x, y) = \frac{xy}{x^2 + y^2}$, $f(0, 0) = 0$. On the axes $f = 0$; on the line $y = x$, $f = \frac{1}{2}$. So $f$ has no limit at $\mathbf{0}$, although $t \mapsto f(t, 0)$ and $t \mapsto f(0, t)$ are continuous.
    \end{example}
    \small
    In several variables, $\bx$ can approach $\mathbf{a}$ from infinitely many directions; checking the coordinate axes is not enough. The same will happen with derivatives: partial derivatives alone do not make a function differentiable.
\end{frame}

\begin{frame}{One variable: review}
    \small
    \begin{proposition}[Linear approximation][prop:linear-approx-1d]
        $f'(x) = c$ $\iff$ $f(x + h) = f(x) + ch + r(h)$ with $r(h)/h \to 0$.
    \end{proposition}
    Sum, product, quotient and chain rule: $(\varphi \circ g)'(x) = \varphi'(g(x))\,g'(x)$.
    \begin{example}[Activation functions][ex:sigmoid]
        \small $\sigma(x) = \frac{1}{1 + e^{-x}}$: \ $\sigma'(x) = \sigma(x)(1 - \sigma(x)) \in (0, \frac{1}{4}]$. \ $\tanh' = 1 - \tanh^2$. \ $\ReLU' = 0$ for $x < 0$ and $1$ for $x > 0$, not differentiable at $0$ (software uses $0$).
    \end{example}
    \sidebyside{0.52}{%
        \begin{theorem}[Taylor][thm:taylor-1d]
            \small For some $\xi$ between $x_0$ and $x$:
            \[f(x) = \sum_{j=0}^k\frac{f^{(j)}(x_0)}{j!}(x - x_0)^j + \frac{f^{(k+1)}(\xi)}{(k+1)!}(x - x_0)^{k+1}.\]
        \end{theorem}
    }{%
        \centering\notesfig[0.8]{L3-taylor-cos}
    }
\end{frame}

\begin{frame}{Worked examples: derivatives and Taylor polynomials}
    \small
    \textbf{Sigmoid.} $\sigma(x) = (1 + e^{-x})^{-1}$; chain rule: $\sigma'(x) = -(1 + e^{-x})^{-2}\cdot(-e^{-x}) = \frac{e^{-x}}{(1 + e^{-x})^2}$. Since $1 - \sigma = \frac{e^{-x}}{1 + e^{-x}}$:
    \[\sigma'(x) = \frac{1}{1 + e^{-x}}\cdot\frac{e^{-x}}{1 + e^{-x}} = \sigma(x)\big(1 - \sigma(x)\big), \qquad \sigma'(0) = \tfrac{1}{2}\cdot\tfrac{1}{2} = \tfrac{1}{4} = \max\sigma'.\]
    \textbf{Taylor for $e^x$ at $0$.} All derivatives are $e^0 = 1$: \ $T_2(x) = 1 + x + \frac{x^2}{2}$, \ $e^{0.1} \approx 1 + 0.1 + 0.005 = 1.105$.
    \\ Error bound: $\frac{e^\xi}{3!}(0.1)^3 \leq \frac{e^{0.1}}{6}\cdot 0.001 \approx 1.8 \cdot 10^{-4}$; \ true value $e^{0.1} = 1.10517\ldots$, error $1.7 \cdot 10^{-4}$ \checkmark.
    \\[0.4em]
    \textbf{Finite differences} for $f = e^x$ at $x = 1$, $h = 0.1$ ($f'(1) = e \approx 2.71828$):
    \[\frac{e^{1.1} - e^{1}}{0.1} \approx \frac{3.00417 - 2.71828}{0.1} = 2.8588, \qquad \text{error } 0.1406 \leq \tfrac{M_2}{2}h = \tfrac{e^{1.1}}{2}\cdot 0.1 \approx 0.150;\]
    \[\frac{e^{1.1} - e^{0.9}}{0.2} \approx \frac{3.00417 - 2.45960}{0.2} = 2.72281, \qquad \text{error } 0.0045 \leq \tfrac{M_3}{6}h^2 \approx \tfrac{3.004}{6}\cdot 0.01 \approx 0.0050.\]
    Halving $h$ halves the first error and quarters the second.
\end{frame}

\begin{frame}{Finite differences}
    \sidebyside{0.48}{%
        \begin{proposition}[][prop:finite-differences]
            \small With $|f''| \leq M_2$, $|f'''| \leq M_3$:
            \[\Big|\tfrac{f(x + h) - f(x)}{h} - f'(x)\Big| \leq \tfrac{M_2}{2}h,\]
            \[\Big|\tfrac{f(x + h) - f(x - h)}{2h} - f'(x)\Big| \leq \tfrac{M_3}{6}h^2.\]
        \end{proposition}
        \small
        Proof: Taylor at $x$ with $k = 1$ and $k = 2$; in the central difference the $f''$ terms cancel. \\[0.3em]
        Too small $h$: rounding errors $\approx \frac{10^{-16}}{h}$.
    }{%
        \centering\fig[1]{s07_finite_differences}
    }
    \begin{mlbox}
        \small Finite differences need $n + 1$ evaluations for a gradient in $\R^n$: hopeless for $n = 10^6$, but perfect for \emph{checking} a hand-written gradient (gradient checking).
    \end{mlbox}
\end{frame}

% ------------------------------------------------------------------ Gradients
\topic{\secref{sec:gradients}}{Partial derivatives and gradients}

\begin{frame}{Partial derivatives and the gradient}
    \begin{definition}[][def:partial]
        \small $\frac{\partial f}{\partial x_i}(\bx) = \lim_{h \to 0}\frac{f(\bx + h\be_i) - f(\bx)}{h}$, and the gradient is the column vector $\nabla f(\bx) = \big(\frac{\partial f}{\partial x_1}, \ldots, \frac{\partial f}{\partial x_n}\big)^\top$.
    \end{definition}
    \begin{example}[][ex:partials]
        \small $f = x_1^2x_2 + x_1x_2^3$: \ $\frac{\partial f}{\partial x_1} = 2x_1x_2 + x_2^3$, $\frac{\partial f}{\partial x_2} = x_1^2 + 3x_1x_2^2$, \ $\nabla f(1, 2) = (12, 13)^\top$.
    \end{example}
    \begin{proposition}[][prop:basic-gradients]
        \small
        \[\nabla(\mathbf{a}^\top\bx) = \mathbf{a}, \qquad \nabla(\bx^\top A\bx) = (A + A^\top)\bx \ \ (= 2A\bx \text{ if } A = A^\top), \qquad \nabla\norm{\bx} = \frac{\bx}{\norm{\bx}}.\]
    \end{proposition}
    \small In ML software, the gradient with respect to a parameter has the shape of the parameter.
\end{frame}

\board{Compute the basic gradients}{$\frac{\partial}{\partial x_k}\sum_{i,j}A_{ij}x_ix_j = ?$ \quad Then $\nabla\norm{\bx}^2 = 2\bx$ and $\nabla\norm{\bx} = \nabla\sqrt{\norm{\bx}^2}$ by the one-variable chain rule.}

\checkslide{the basic gradients}{
    \small
    \textbf{Quadratic form.} In $\sum_{i,j}A_{ij}x_ix_j$, the variable $x_k$ appears when $i = k$ or $j = k$ (product rule for $i = j = k$):
    \[\frac{\partial}{\partial x_k}\sum_{i,j}A_{ij}x_ix_j = \sum_jA_{kj}x_j + \sum_iA_{ik}x_i = (A\bx)_k + (A^\top\bx)_k.\]
    \textbf{Numbers.} $A = \begin{pmatrix} 1 & 2 \\ 0 & 3 \end{pmatrix}$: $f(\bx) = x_1^2 + 2x_1x_2 + 3x_2^2$. Directly: $\nabla f = (2x_1 + 2x_2,\ 2x_1 + 6x_2)^\top$; at $\bx = (1, 1)^\top$: $(4, 8)^\top$. \ Formula: $(A + A^\top)\bx = \begin{pmatrix} 2 & 2 \\ 2 & 6 \end{pmatrix}\begin{pmatrix} 1 \\ 1 \end{pmatrix} = \begin{pmatrix} 4 \\ 8 \end{pmatrix}$ \checkmark.
    \\[0.4em]
    \textbf{Norms.} $A = I$: $\nabla\norm{\bx}^2 = 2\bx$. \ $\norm{\bx} = g(\norm{\bx}^2)$ with $g(s) = \sqrt s$, $g'(s) = \frac{1}{2\sqrt s}$:
    \[\nabla\norm{\bx} = \frac{1}{2\norm{\bx}}\cdot 2\bx = \frac{\bx}{\norm{\bx}} \quad (\bx \neq \mathbf{0}), \qquad \text{e.g. } \bx = (3, 4)^\top: \ \nabla\norm{\bx} = (0.6, 0.8)^\top,\]
    a unit vector: the norm grows fastest straight away from the origin.
    \\[0.3em]
    \textbf{Linear.} $\mathbf{a}^\top\bx = \sum_ia_ix_i$: $\frac{\partial}{\partial x_k} = a_k$, so $\nabla(\mathbf{a}^\top\bx) = \mathbf{a}$.
}

\begin{frame}{Differentiability: the derivative is a linear map}
    \begin{definition}[][def:differentiable]
        \small $f: D \to \R$ is differentiable at $\bx$ if there is a row vector $\mathbf{c} \in \R^{1 \times n}$ with
        \[f(\bx + \mathbf{h}) = f(\bx) + \mathbf{c}\mathbf{h} + o(\norm{\mathbf{h}}).\]
        Then $Df(\bx) = \mathbf{c}$ is the derivative.
    \end{definition}
    \begin{proposition}[][prop:differentiable]
        \small If $f$ is differentiable at $\bx$: every directional derivative exists and $D_{\bv}f(\bx) = \nabla f(\bx)^\top\bv$; $Df(\bx) = \nabla f(\bx)^\top$; $f$ is continuous at $\bx$.
    \end{proposition}
    \begin{theorem}[][thm:C1]
        \small Partial derivatives that exist near $\bx$ and are continuous at $\bx$ $\Rightarrow$ $f$ is differentiable at $\bx$.
    \end{theorem}
    \small The function $\frac{xy}{x^2 + y^2}$ has both partial derivatives at $\mathbf{0}$ (they are $0$), but it is not even continuous there.
\end{frame}

\begin{frame}{Derivative vs gradient}
    \begin{remark}[The gradient depends on the inner product][rem:riesz]
        \small $Df(\bx)$ is a linear functional, an element of $(\R^n)^*$ (Lecture~1). The gradient is the vector that represents it with the inner product (Riesz, Lecture~6):
        \[Df(\bx)\mathbf{h} = \ip{\nabla f(\bx)}{\mathbf{h}} \quad \text{for all } \mathbf{h}.\]
        For $\ip{\bx}{\by}_A = \bx^\top A\by$ with SPD $A$: \ $Df(\bx)\mathbf{h} = \nabla f(\bx)^\top\mathbf{h} = (A^{-1}\nabla f(\bx))^\top A\,\mathbf{h}$, so the gradient becomes
        \[\nabla_A f(\bx) = A^{-1}\nabla f(\bx).\]
    \end{remark}
    \small
    \begin{itemize}
        \item The derivative depends only on $f$; the gradient also on how we measure lengths.
        \item Changing the inner product $=$ preconditioning. Newton's method ($A = $ Hessian, Lecture~8), the natural gradient ($A = $ Fisher information), Adam (diagonal $A$).
    \end{itemize}
\end{frame}

\begin{frame}{Steepest ascent}
    \begin{theorem}[Steepest ascent][thm:steepest]
        \small If $\nabla f(\bx) \neq \mathbf{0}$, then for every unit vector $\bv$:
        \(-\norm{\nabla f(\bx)} \leq D_{\bv}f(\bx) \leq \norm{\nabla f(\bx)},\)
        with the maximum only at $\bv = \nabla f/\norm{\nabla f}$ and the minimum only at $-\nabla f/\norm{\nabla f}$.
    \end{theorem}
    \centering\fig[0.66]{s07_landscape}
    \\[0.1em]
    \raggedright\small $\nabla f(1, 2) = (12, 13)^\top$: in the direction $\frac{1}{5}(3, 4)^\top$ the rate is $17.6$; the maximum is $\sqrt{313} \approx 17.69$.
\end{frame}

\board{Proof of steepest ascent}{$D_{\bv}f(\bx) = \nabla f(\bx)^\top\bv$ and Cauchy--Schwarz. When is there equality?}

\checkslide{directional derivatives with numbers}{
    \small
    $f(x_1, x_2) = x_1^2x_2 + x_1x_2^3$: \ $\frac{\partial f}{\partial x_1} = 2x_1x_2 + x_2^3$, \ $\frac{\partial f}{\partial x_2} = x_1^2 + 3x_1x_2^2$; at $(1, 2)$:
    \[\nabla f(1, 2) = (2 \cdot 1 \cdot 2 + 8,\ 1 + 3 \cdot 1 \cdot 4)^\top = (12, 13)^\top, \qquad \norm{\nabla f} = \sqrt{144 + 169} = \sqrt{313} \approx 17.69.\]
    \begin{center}
        \renewcommand{\arraystretch}{1.3}
        \begin{tabular}{@{}lll@{}}
            \toprule
            unit direction $\bv$ & $D_{\bv}f = \nabla f^\top\bv$ & \\
            \midrule
            $\be_1$ & $12$ & the partial derivative $\frac{\partial f}{\partial x_1}$ \\
            $\frac{1}{5}(3, 4)^\top$ & $\frac{1}{5}(36 + 52) = 17.6$ & close to the maximum \\
            $\frac{1}{\sqrt{313}}(12, 13)^\top$ & $\frac{144 + 169}{\sqrt{313}} = \sqrt{313} \approx 17.69$ & steepest ascent \\
            $\frac{1}{\sqrt{313}}(13, -12)^\top$ & $\frac{156 - 156}{\sqrt{313}} = 0$ & along the level curve \\
            $-\frac{1}{\sqrt{313}}(12, 13)^\top$ & $-17.69$ & steepest descent \\
            \bottomrule
        \end{tabular}
    \end{center}
    Check against the definition with $t = 0.005$, $\bv = (0.6, 0.8)$: $\frac{f(1.003, 2.004) - f(1, 2)}{0.005} \approx \frac{10.0883 - 10}{0.005} \approx 17.66$; the gap to $17.6$ is the $O(t)$ error.
}

% ------------------------------------------------------------------ Jacobians
\topic{\secref{sec:jacobians}}{Jacobians}

\begin{frame}{The Jacobian matrix}
    \begin{definition}[][def:jacobian]
        \small $f: D \to \R^m$ is differentiable at $\bx$ if $f(\bx + \mathbf{h}) = f(\bx) + J\mathbf{h} + o(\norm{\mathbf{h}})$ for an $m \times n$ matrix $J = Df(\bx)$.
    \end{definition}
    \begin{proposition}[][prop:jacobian-entries]
        \small $\big(Df(\bx)\big)_{ij} = \frac{\partial f_i}{\partial x_j}(\bx)$: \ row $i$ $=$ $\nabla f_i(\bx)^\top$ (outputs), column $j$ $=$ $\frac{\partial f}{\partial x_j}$ (inputs).
    \end{proposition}
    \small
    \begin{example}[][ex:jacobian]
        $f(x_1, x_2) = (x_1^2x_2, \, 3x_1 + x_2^2)^\top$: \ $Df(\bx) = \begin{pmatrix} \partial_1f_1 & \partial_2f_1 \\ \partial_1f_2 & \partial_2f_2 \end{pmatrix} = \begin{pmatrix} 2x_1x_2 & x_1^2 \\ 3 & 2x_2 \end{pmatrix}$. At $\mathbf{a} = (1, 2)^\top$: $f(\mathbf{a}) = (1 \cdot 2,\ 3 + 4)^\top = (2, 7)^\top$, $Df(\mathbf{a}) = \begin{pmatrix} 4 & 1 \\ 3 & 4 \end{pmatrix}$, and for $\mathbf{h} = (0.1, -0.1)^\top$:
        \[f(1.1, 1.9) \approx \begin{pmatrix} 2 \\ 7 \end{pmatrix} + \begin{pmatrix} 0.4 - 0.1 \\ 0.3 - 0.4 \end{pmatrix} = \begin{pmatrix} 2.3 \\ 6.9 \end{pmatrix}; \quad \text{exact: } \begin{pmatrix} 1.21 \cdot 1.9 \\ 3.3 + 3.61 \end{pmatrix} = \begin{pmatrix} 2.299 \\ 6.91 \end{pmatrix}.\]
    \end{example}
    \begin{itemize}
        \item affine $f(\bx) = A\bx + \mathbf{b}$: \ $Df = A$;
        \item entrywise $\varphi$: \ $Df(\bx) = \diag(\varphi'(x_1), \ldots, \varphi'(x_n))$, \ $Df(\bx)\bv = \varphi'(\bx) \odot \bv$.
    \end{itemize}
\end{frame}

\board{Example: the Jacobian of softmax}{$s_i = e^{z_i}/S$, $S = \sum_j e^{z_j}$. Show $\frac{\partial s_i}{\partial z_j} = s_i\delta_{ij} - s_is_j$, i.e., $D\softmax(\mathbf{z}) = \diag(\mathbf{s}) - \mathbf{s}\mathbf{s}^\top$. Why do the rows sum to $0$?}

\checkslide{the softmax Jacobian}{
    \small
    \textbf{Derivation} (quotient rule, $\frac{\partial S}{\partial z_j} = e^{z_j}$): \ $\frac{\partial s_i}{\partial z_j} = \frac{\delta_{ij}e^{z_i}S - e^{z_i}e^{z_j}}{S^2} = s_i\delta_{ij} - s_is_j$.
    \\[0.3em]
    \textbf{Numbers.} $\mathbf{z} = (2, 1, 0)^\top$: $e^2 \approx 7.389$, $e^1 \approx 2.718$, $e^0 = 1$, $S \approx 11.107$, \ $\mathbf{s} \approx (0.665, 0.245, 0.090)^\top$. Entries: $s_1 - s_1^2 \approx 0.665 - 0.442 = 0.223$, \ $-s_1s_2 \approx -0.163$, \ $s_3 - s_3^2 \approx 0.090 - 0.008 = 0.082$, \ldots
    \[D\softmax(\mathbf{z}) = \diag(\mathbf{s}) - \mathbf{s}\mathbf{s}^\top \approx \begin{pmatrix} 0.223 & -0.163 & -0.060 \\ -0.163 & 0.185 & -0.022 \\ -0.060 & -0.022 & 0.082 \end{pmatrix}, \ \ \text{row sum } 0.223 - 0.163 - 0.060 = 0 \,\checkmark.\]
    \begin{itemize}
        \item Symmetric, and $D\softmax(\mathbf{z})\mathbf{1} = \mathbf{s} - \mathbf{s}(\mathbf{s}^\top\mathbf{1}) = \mathbf{0}$.
        \item Reason: $\softmax(\mathbf{z} + c\mathbf{1}) = \softmax(\mathbf{z})$; moving along $\mathbf{1}$ changes nothing. \\ (Lecture~1: softmax lives on the quotient space $\R^K/\Span\{\mathbf{1}\}$.)
        \item In practice one computes $\softmax(\mathbf{z} - \max_i z_i)$ to avoid overflow.
    \end{itemize}
}

\begin{frame}{The Jacobian determinant}
    \sidebyside{0.5}{%
        \begin{example}[Polar coordinates][ex:polar]
            \small $f(r, \varphi) = (r\cos\varphi, r\sin\varphi)^\top$,
            \[Df = \begin{pmatrix} \cos\varphi & -r\sin\varphi \\ \sin\varphi & r\cos\varphi \end{pmatrix}, \ \det Df = r.\]
        \end{example}
        \small $|\det Df(\bx)|$ is the local factor by which $f$ scales areas and volumes:
        \[\textstyle\iint g\,dx\,dy = \int_0^{2\pi}\!\!\int_0^\infty g\,r\,dr\,d\varphi.\]
    }{%
        \centering\notesfig[1]{L3-polar}
    }
    \small
    Gaussian integral: $\big(\int e^{-x^2}dx\big)^2 = \int_0^{2\pi}\!\int_0^\infty e^{-r^2}r\,dr\,d\varphi = \pi$, so $\int_{-\infty}^\infty e^{-x^2}dx = \sqrt{\pi}$.\\
    \textcolor{mlcolor}{Normalizing flows transform densities with $|\det Df|$; they need cheap determinants.}
\end{frame}

% ------------------------------------------------------------------ Chain rule
\topic{\secref{sec:chain-rule}}{The chain rule}

\begin{frame}{The chain rule}
    \begin{theorem}[Chain rule][thm:chain-rule]
        \small If $g$ is differentiable at $\bx$ and $f$ at $g(\bx)$, then
        \[D(f \circ g)(\bx) = \underbrace{Df(g(\bx))}_{k \times m}\,\underbrace{Dg(\bx)}_{m \times n}, \qquad \frac{\partial(f \circ g)_i}{\partial x_j} = \sum_{l=1}^m\frac{\partial f_i}{\partial y_l}\frac{\partial g_l}{\partial x_j}.\]
    \end{theorem}
    \begin{corollary}[][cor:chain-special]
        \small
        \begin{enumerate}
            \item Along a curve: $\frac{d}{dt}f(\bx(t)) = \nabla f(\bx(t))^\top\bx'(t)$.
            \item Scalar $f$: \ $\nabla(f \circ g)(\bx) = Dg(\bx)^\top\,\nabla f(g(\bx))$ \ (gradients are pulled back by the transposed Jacobian).
        \end{enumerate}
    \end{corollary}
    \small Composition of linear approximations $=$ product of matrices. Write the sizes under the factors: it checks the order and the transposes.
\end{frame}

\begin{frame}{Worked examples: the chain rule two ways}
    \small
    \textbf{Along a curve.} $f(x_1, x_2) = x_1^2 + 2x_2$, \ $\bx(t) = (\sin t, \cos t)^\top$ (the unit circle):
    \[\tfrac{d}{dt}f(\bx(t)) = \nabla f^\top\bx'(t) = \begin{pmatrix} 2\sin t \\ 2 \end{pmatrix}^\top\begin{pmatrix} \cos t \\ -\sin t \end{pmatrix} = 2\sin t\cos t - 2\sin t.\]
    Directly: $f(\bx(t)) = \sin^2t + 2\cos t$, derivative $2\sin t\cos t - 2\sin t$ \checkmark. \ At $t = \frac{\pi}{2}$: $0 - 2 = -2$.
    \\[0.5em]
    \textbf{A composition.} $g(x, y) = (xy, \ x + y)^\top$, $f(u, v) = u^2 + v$; gradient of $f \circ g$ at $(1, 2)$:
    \[g(1, 2) = (2, 3)^\top, \quad \nabla f(g) = (2u, 1)^\top = (4, 1)^\top, \quad Dg = \begin{pmatrix} y & x \\ 1 & 1 \end{pmatrix} = \begin{pmatrix} 2 & 1 \\ 1 & 1 \end{pmatrix},\]
    \[\nabla(f \circ g) = \underbrace{Dg^\top}_{2 \times 2}\underbrace{\nabla f(g)}_{2 \times 1} = \begin{pmatrix} 2 & 1 \\ 1 & 1 \end{pmatrix}\begin{pmatrix} 4 \\ 1 \end{pmatrix} = \begin{pmatrix} 9 \\ 5 \end{pmatrix}.\]
    Directly: $(f \circ g)(x, y) = x^2y^2 + x + y$, \ $\partial_x = 2xy^2 + 1 = 9$, \ $\partial_y = 2x^2y + 1 = 5$ \checkmark.
\end{frame}

\begin{frame}{The gradient is orthogonal to level sets}
    \sidebyside{0.5}{%
        \begin{corollary}[][cor:level-sets]
            \small If $f(\bx(t)) = c$ for all $t$, then $\nabla f(\bx(t))^\top\bx'(t) = 0$.
        \end{corollary}
        \small Differentiate $f(\bx(t)) = c$ with the chain rule. \\[0.4em]
        Level curves of $f = x_1^2 + 4x_2^2$ and $\nabla f = (2x_1, 8x_2)^\top$ (scaled by $0.1$) on the level curve $c = 9$. \\[0.4em]
        Gradient descent moves orthogonally to the level sets, not towards the minimum.
    }{%
        \centering\notesfig[1]{L3-gradient-field}
    }
\end{frame}

\board{Example: the gradient of least squares}{$L(\boldsymbol{\theta}) = \norm{\by - \Phi\boldsymbol{\theta}}^2 = f(\mathbf{e}(\boldsymbol{\theta}))$ with $\mathbf{e} = \by - \Phi\boldsymbol{\theta}$, $f(\mathbf{e}) = \norm{\mathbf{e}}^2$. Chain rule with sizes; set $\nabla L = \mathbf{0}$.}

\board{Example: logistic regression}{$\ell(\bw) = -y\log p - (1 - y)\log(1 - p)$, \ $p = \sigma(z)$, \ $z = \bw^\top\bx$. Show $\nabla\ell(\bw) = (p - y)\bx$. \\ Numbers: $\bx = (1, 2)^\top$, $y = 1$, $\bw = (0.5, -0.25)^\top$.}

\checkslide{least squares and logistic regression}{
    \small
    Least squares: $DL(\boldsymbol{\theta}) = \underbrace{2(\by - \Phi\boldsymbol{\theta})^\top}_{1 \times N}\underbrace{(-\Phi)}_{N \times d}$, so
    \[\nabla L(\boldsymbol{\theta}) = 2\Phi^\top(\Phi\boldsymbol{\theta} - \by) = \mathbf{0} \iff \Phi^\top\Phi\boldsymbol{\theta} = \Phi^\top\by \quad \text{(the normal equations of Lecture~6).}\]
    Logistic regression, along $\bw \mapsto z \mapsto p \mapsto \ell$: \ $\frac{d\ell}{dp} = -\frac{y}{p} + \frac{1 - y}{1 - p}$, \ $\frac{dp}{dz} = p(1 - p)$, \ $Dz(\bw) = \bx^\top$:
    \[\tfrac{d\ell}{dz} = -y(1 - p) + (1 - y)p = p - y, \qquad \nabla\ell(\bw) = \big(\sigma(\bw^\top\bx) - y\big)\,\bx \quad \text{(prediction error $\times$ input).}\]
    \textbf{Numbers:} $z = 0.5 \cdot 1 - 0.25 \cdot 2 = 0$, $p = \sigma(0) = \frac{1}{2}$, $\ell = -\log\frac{1}{2} = \log 2 \approx 0.693$, $\nabla\ell = (\frac{1}{2} - 1)(1, 2)^\top = (-0.5, -1)^\top$.
    \\ One gradient step with $\eta = 1$: $\bw = (1, 0.75)^\top$, $z = 2.5$, $p \approx 0.924$, $\ell \approx 0.079$: the loss dropped.
    \begin{mlbox}
        \small Softmax with cross-entropy: $\nabla_{\mathbf{z}}(-\log p_c) = \mathbf{p} - \be_c$. Same form: ``prediction minus target''.
    \end{mlbox}
}

% ------------------------------------------------------------------ Matrix gradients
\topic{\secref{sec:matrix-gradients}}{Gradients with respect to matrices}

\begin{frame}{Gradients with respect to matrices}
    \small
    $f: \R^{m \times n} \to \R$: \ $\big(\nabla f(X)\big)_{ij} = \frac{\partial f}{\partial X_{ij}}$, a matrix of the same shape as $X$.
    \\[0.3em]
    With the Frobenius inner product $\ip{A}{B} = \tr(A^\top B)$:
    \[f(X + H) = f(X) + \ip{\nabla f(X)}{H} + o(\norm{H}_F).\]
    \begin{block}{The method}
        Expand $f(X + H)$, keep the terms linear in $H$, and write them as $\tr(G^\top H)$. Then $\nabla f(X) = G$. \\
        Tools: $\tr(A) = \tr(A^\top)$, \ $\tr(AB) = \tr(BA)$, \ hence $\tr(MH) = \ip{M^\top}{H}$.
    \end{block}
    \begin{proposition}[][prop:matrix-gradients]
        \[\nabla_X\tr(AX) = A^\top, \qquad \nabla_X(\mathbf{a}^\top X\mathbf{b}) = \mathbf{a}\mathbf{b}^\top,\]
        \[\nabla_X\norm{AX - B}_F^2 = 2A^\top(AX - B), \qquad \nabla_X\tr(X^\top AX) = (A + A^\top)X.\]
    \end{proposition}
\end{frame}

\board[2]{Use the method}{(1) $\norm{A(X + H) - B}_F^2 = \norm{R + AH}_F^2$ with $R = AX - B$. \quad (2) $\tr((X + H)^\top A(X + H))$. \quad (3) $\mathbf{a}^\top X\mathbf{b} = \tr(\mathbf{b}\mathbf{a}^\top X)$.}

\checkslide{gradients with respect to matrices}{
    \small
    \textbf{(1)} $\norm{R + AH}_F^2 = \norm{R}_F^2 + 2\ip{R}{AH} + \norm{AH}_F^2$; the last term is $O(\norm{H}^2)$, and
    \[\ip{R}{AH} = \tr(R^\top AH) = \tr\big((A^\top R)^\top H\big) = \ip{A^\top R}{H} \ \Rightarrow \ \nabla_X\norm{AX - B}_F^2 = 2A^\top(AX - B).\]
    \textbf{(2)} $\tr((X + H)^\top A(X + H)) = \tr(X^\top AX) + \tr(H^\top AX) + \tr(X^\top AH) + O(\norm{H}^2)$, with
    \[\tr(H^\top AX) = \ip{AX}{H}, \qquad \tr(X^\top AH) = \ip{A^\top X}{H} \ \Rightarrow \ \nabla_X = (A + A^\top)X.\]
    \textbf{(3)} $\mathbf{a}^\top X\mathbf{b} = \sum_{i,j}a_iX_{ij}b_j$, so $\frac{\partial}{\partial X_{ij}} = a_ib_j$: \ $\nabla_X = \mathbf{a}\mathbf{b}^\top$. Numbers: $\mathbf{a} = (1, 2)^\top$, $\mathbf{b} = (3, 1)^\top$:
    \[\mathbf{a}^\top X\mathbf{b} = 3X_{11} + X_{12} + 6X_{21} + 2X_{22}, \qquad \nabla_X = \begin{pmatrix} 3 & 1 \\ 6 & 2 \end{pmatrix} = \begin{pmatrix} 1 \\ 2 \end{pmatrix}\begin{pmatrix} 3 & 1 \end{pmatrix} \ \checkmark.\]
    \textbf{Jacobi, $2 \times 2$.} $X = \begin{pmatrix} a & b \\ c & d \end{pmatrix}$, $\det X = ad - bc$: \ $\nabla\det = \begin{pmatrix} d & -c \\ -b & a \end{pmatrix}$, and $X^{-1} = \frac{1}{ad - bc}\begin{pmatrix} d & -b \\ -c & a \end{pmatrix}$, so $\det(X)X^{-\top} = \begin{pmatrix} d & -c \\ -b & a \end{pmatrix}$ \checkmark.
}

\begin{frame}{Jacobi's formula}
    \begin{lemma}[][lem:det-identity]
        \small $\det(I + tE) = 1 + t\tr(E) + t^2q(t)$ for a polynomial $q$; so $\frac{d}{dt}\det(I + tE)\big|_{t=0} = \tr(E)$.
    \end{lemma}
    \begin{proposition}[Jacobi's formula][prop:det-gradient]
        \small For invertible $X$: \quad $\nabla_X\det(X) = \det(X)\,X^{-\top}$, \qquad $\nabla_X\log|\det X| = X^{-\top}$.
    \end{proposition}
    \small
    Proof: $\det(X + tH) = \det(X)\det(I + tX^{-1}H) = \det(X)\big(1 + t\tr(X^{-1}H) + O(t^2)\big)$, and $\det(X)\tr(X^{-1}H) = \ip{\det(X)X^{-\top}}{H}$.
    \\[0.3em]
    Check: $n = 1$ gives $\frac{d}{dx}\log|x| = \frac{1}{x}$. Also: $\frac{\partial\det X}{\partial X_{ij}}$ is the cofactor $C_{ij}$, and $\det(X)X^{-\top} = C$ (Lecture~1).
    \begin{mlbox}
        \small $\log|\det|$ appears in Gaussian log-likelihoods ($\log\det\Sigma$) and in normalizing flows.
    \end{mlbox}
\end{frame}

\begin{frame}{A linear layer}
    \begin{example}[][ex:linear-layer]
        \small $\by = W\bx + \mathbf{b}$, a loss $L$ depending on $\by$, and $\boldsymbol{\delta} = \nabla_{\by}L$. Then
        \[\nabla_W L = \boldsymbol{\delta}\,\bx^\top, \qquad \nabla_{\mathbf{b}}L = \boldsymbol{\delta}, \qquad \nabla_{\bx}L = W^\top\boldsymbol{\delta}.\]
    \end{example}
    \small
    \begin{itemize}
        \item $\frac{\partial L}{\partial W_{ij}} = \frac{\partial L}{\partial y_i}\frac{\partial y_i}{\partial W_{ij}} = \delta_ix_j$, an outer product; the sizes fit: $(m \times 1)(1 \times n)$.
        \item A batch: inputs are the rows of $X$ ($N \times n$), $Y = XW^\top + \mathbf{1}\mathbf{b}^\top$; \ $\nabla_W L = \sum_k\boldsymbol{\delta}_k\bx_k^\top = \Delta^\top X$, \ $\nabla_{\mathbf{b}}L = \Delta^\top\mathbf{1}$.
        \item Softmax $+$ cross-entropy on top: $\boldsymbol{\delta} = \mathbf{p} - \be_c$, so $\nabla_W\ell = (\mathbf{p} - \be_c)\bx^\top$.
    \end{itemize}
\end{frame}

\begin{frame}{Worked example: a linear layer and softmax with numbers}
    \small
    \textbf{Linear layer} $\by = W\bx$ with $W = \begin{pmatrix} 1 & 2 \\ 0 & 1 \end{pmatrix}$, $\bx = (1, 1)^\top$, loss $L = \frac{1}{2}\norm{\by - \mathbf{t}}^2$, target $\mathbf{t} = (1, 1)^\top$:
    \[\by = (3, 1)^\top, \qquad \boldsymbol{\delta} = \nabla_{\by}L = \by - \mathbf{t} = (2, 0)^\top,\]
    \[\nabla_WL = \boldsymbol{\delta}\bx^\top = \begin{pmatrix} 2 \\ 0 \end{pmatrix}\begin{pmatrix} 1 & 1 \end{pmatrix} = \begin{pmatrix} 2 & 2 \\ 0 & 0 \end{pmatrix}, \qquad \nabla_{\bx}L = W^\top\boldsymbol{\delta} = \begin{pmatrix} 1 & 0 \\ 2 & 1 \end{pmatrix}\begin{pmatrix} 2 \\ 0 \end{pmatrix} = \begin{pmatrix} 2 \\ 4 \end{pmatrix}.\]
    Check one entry: $L = \frac{1}{2}\big((W_{11}x_1 + W_{12}x_2 - 1)^2 + (W_{21}x_1 + W_{22}x_2 - 1)^2\big)$, so $\frac{\partial L}{\partial W_{11}} = (3 - 1) \cdot x_1 = 2$ \checkmark. The second row gets no gradient: $y_2$ is already correct.
    \\[0.5em]
    \textbf{Softmax with cross-entropy}, $\mathbf{z} = (2, 1, 0)^\top$, correct class $c = 1$:
    \[\mathbf{p} \approx (0.665, 0.245, 0.090)^\top, \qquad \ell = -\log p_1 \approx 0.408, \qquad \nabla_{\mathbf{z}}\ell = \mathbf{p} - \be_1 \approx (-0.335, 0.245, 0.090)^\top.\]
    A gradient step raises $z_1$ and lowers $z_2$, $z_3$; the gradient sums to $0$ (shifting all scores changes nothing). If $\mathbf{z} = W\bx + \mathbf{b}$: \ $\nabla_W\ell = (\mathbf{p} - \be_1)\bx^\top$.
\end{frame}

\begin{frame}{Gradients to remember}
    \centering\small
    \renewcommand{\arraystretch}{1.3}
    \begin{tabular}{@{}lll@{}}
        \toprule
        function & gradient & assumptions \\
        \midrule
        $\mathbf{a}^\top\bx$ & $\mathbf{a}$ & \\
        $\bx^\top A\bx$ & $(A + A^\top)\bx$ & $= 2A\bx$ if $A = A^\top$ \\
        $\norm{\bx}^2$, \ $\norm{\bx}$ & $2\bx$, \ $\bx/\norm{\bx}$ & $\bx \neq \mathbf{0}$ \\
        $\norm{A\bx - \mathbf{b}}^2$ & $2A^\top(A\bx - \mathbf{b})$ & \\
        $f(g(\bx))$ & $Dg(\bx)^\top\nabla f(g(\bx))$ & $f$ scalar \\
        $\log(e^{z_1} + \cdots + e^{z_K})$ & $\softmax(\mathbf{z})$ & \\
        $-\log\softmax(\mathbf{z})_c$ & $\softmax(\mathbf{z}) - \be_c$ & \\
        $\tr(AX)$, \ $\mathbf{a}^\top X\mathbf{b}$ & $A^\top$, \ $\mathbf{a}\mathbf{b}^\top$ & \\
        $\norm{AX - B}_F^2$ & $2A^\top(AX - B)$ & \\
        $\tr(X^\top AX)$ & $(A + A^\top)X$ & \\
        $\det X$, \ $\log|\det X|$ & $\det(X)X^{-\top}$, \ $X^{-\top}$ & $X$ invertible \\
        \bottomrule
    \end{tabular}
\end{frame}

% ------------------------------------------------------------------ Autodiff
\topic{\secref{sec:autodiff}}{Automatic differentiation and backpropagation}

\begin{frame}{Three ways to compute derivatives}
    \small
    \textbf{Symbolic}: formulas explode, and programs with loops are not one formula. \textbf{Numerical}: $n + 1$ evaluations, approximate. \textbf{Automatic} (AD): apply the chain rule to the elementary operations of the program ($+$, $\times$, $\exp$, $\log$, \dots), with \emph{numbers}; exact up to rounding.
    \\[0.4em]
    For $f = f_L \circ \cdots \circ f_1$ with $\bx_l = f_l(\bx_{l-1})$:
    \[Df(\bx) = J_LJ_{L-1}\cdots J_1, \qquad J_l = Df_l(\bx_{l-1}) \in \R^{n_l \times n_{l-1}}.\]
    \begin{itemize}
        \item \textbf{Forward mode}: $J_L(\cdots(J_2(J_1\bv)))$, a Jacobian-vector product; one column per pass.
        \item \textbf{Reverse mode}: $((\bu^\top J_L)J_{L-1})\cdots J_1$, a vector-Jacobian product; one row per pass; stores the $\bx_l$ (memory).
    \end{itemize}
    For a loss ($m = 1$), one reverse pass gives the whole gradient, for $2$--$5$ evaluations of $f$, whatever $n$ is. Reverse mode applies the dual maps: $(S \circ T)^* = T^* \circ S^*$.
\end{frame}

\begin{frame}{Computation graphs}
    \sidebyside{0.48}{%
        \small
        $f(x_1, x_2) = \log x_1 + x_1x_2 - \sin x_2$. \\[0.5em]
        \textbf{Forward mode} (tangents), in the order of evaluation:
        \[\dot v_j = \sum_{i:\, v_i \to v_j}\frac{\partial v_j}{\partial v_i}\,\dot v_i.\]
        \textbf{Reverse mode} (adjoints $\bar v_i = \frac{\partial f}{\partial v_i}$), backwards from $\bar f = 1$:
        \[\bar v_i = \sum_{j:\, v_i \to v_j}\bar v_j\,\frac{\partial v_j}{\partial v_i}.\]
        A variable used several times collects several contributions: the sum in the chain rule.
    }{%
        \centering\notesfig[1]{L3-computation-graph}
    }
\end{frame}

\board[2]{Example: both modes by hand}{$f(x_1, x_2) = \log x_1 + x_1x_2 - \sin x_2$ at $(x_1, x_2) = (2, 5)$. \ Evaluate $v_1, \ldots, v_5$; forward mode with $\dot\bx = \be_1$; reverse mode from $\bar v_5 = 1$. ($\sin 5 \approx -0.959$, $\cos 5 \approx 0.284$.)}

\checkslide{automatic differentiation at $(2, 5)$}{
    \small
    \begin{columns}[T]
        \begin{column}{0.56\linewidth}
            \renewcommand{\arraystretch}{1.15}
            \begin{tabular}{@{}lll@{}}
                \toprule
                evaluation & & forward, $\dot\bx = \be_1$ \\
                \midrule
                $v_1 = \log x_1$ & $0.693$ & $\dot v_1 = \dot x_1/x_1 = 0.5$ \\
                $v_2 = x_1x_2$ & $10$ & $\dot v_2 = \dot x_1x_2 + x_1\dot x_2 = 5$ \\
                $v_3 = \sin x_2$ & $-0.959$ & $\dot v_3 = \cos(x_2)\dot x_2 = 0$ \\
                $v_4 = v_1 + v_2$ & $10.693$ & $\dot v_4 = 5.5$ \\
                $v_5 = v_4 - v_3$ & $11.652$ & $\dot v_5 = 5.5$ \\
                \bottomrule
            \end{tabular}
        \end{column}
        \begin{column}{0.4\linewidth}
            \renewcommand{\arraystretch}{1.15}
            \begin{tabular}{@{}l@{}}
                \toprule
                reverse \\
                \midrule
                $\bar v_5 = 1$, \ $\bar v_4 = 1$, \ $\bar v_3 = -1$ \\
                $\bar v_2 = 1$, \ $\bar v_1 = 1$ \\
                $\bar x_2 = \bar v_2x_1 + \bar v_3\cos x_2 = 1.716$ \\
                $\bar x_1 = \bar v_1/x_1 + \bar v_2x_2 = 5.5$ \\
                \bottomrule
            \end{tabular}
        \end{column}
    \end{columns}
    \vskip0.8em
    Forward mode gave one partial derivative ($\frac{\partial f}{\partial x_1} = 5.5$); reverse mode gave the whole gradient $(5.5, 1.716)^\top$ in one pass.
}

\begin{frame}{Backpropagation in a neural network}
    \small
    One hidden layer, squared loss:
    \[\mathbf{z} = W_1\bx + \mathbf{b}_1, \quad \mathbf{a} = \sigma(\mathbf{z}), \quad \hat\by = W_2\mathbf{a} + \mathbf{b}_2, \quad L = \tfrac{1}{2}\norm{\hat\by - \by}^2.\]
    Forward pass: compute and \emph{store} $\mathbf{z}, \mathbf{a}, \hat\by, L$. Backward pass:
    \[\begin{aligned}
        \boldsymbol{\delta}_2 &= \nabla_{\hat\by}L = \hat\by - \by, & \nabla_{W_2}L &= \boldsymbol{\delta}_2\mathbf{a}^\top, \quad \nabla_{\mathbf{b}_2}L = \boldsymbol{\delta}_2, \\
        \boldsymbol{\delta}_1 &= \nabla_{\mathbf{z}}L = \sigma'(\mathbf{z}) \odot (W_2^\top\boldsymbol{\delta}_2), & \nabla_{W_1}L &= \boldsymbol{\delta}_1\bx^\top, \quad \nabla_{\mathbf{b}_1}L = \boldsymbol{\delta}_1.
    \end{aligned}\]
    Deep network $\mathbf{z}_l = W_l\mathbf{a}_{l-1} + \mathbf{b}_l$, $\mathbf{a}_l = \sigma(\mathbf{z}_l)$: the error signals go backwards,
    \[\boldsymbol{\delta}_l = \sigma'(\mathbf{z}_l) \odot \big(W_{l+1}^\top\boldsymbol{\delta}_{l+1}\big), \qquad \nabla_{W_l}L = \boldsymbol{\delta}_l\mathbf{a}_{l-1}^\top, \qquad \nabla_{\mathbf{b}_l}L = \boldsymbol{\delta}_l.\]
    The backward pass of a layer costs about twice its forward pass. PyTorch and JAX record the graph in the forward pass; each operation knows its vector-Jacobian product.
\end{frame}

\board[2]{Derive backpropagation for one hidden layer}{Use the linear layer ($\nabla_W L = \boldsymbol{\delta}\bx^\top$, $\nabla_{\bx}L = W^\top\boldsymbol{\delta}$) and the entrywise Jacobian $D\sigma(\mathbf{z}) = \diag(\sigma'(\mathbf{z}))$. Check all the sizes. \\ Numbers (one unit per layer): $x = 1$, $y = 1$, $w_1 = 0.5$, $b_1 = 0$, $w_2 = 2$, $b_2 = 0$, sigmoid.}

\checkslide{backpropagation with numbers}{
    \small
    \textbf{Forward pass} (store everything):
    \[z = w_1x + b_1 = 0.5, \quad a = \sigma(0.5) \approx 0.6225, \quad \hat y = w_2a + b_2 \approx 1.2449, \quad L = \tfrac{1}{2}(\hat y - y)^2 \approx 0.0300.\]
    \textbf{Backward pass:}
    \[\begin{aligned}
        \delta_2 &= \hat y - y \approx 0.2449, \qquad \tfrac{\partial L}{\partial w_2} = \delta_2a \approx 0.1525, \qquad \tfrac{\partial L}{\partial b_2} = \delta_2 \approx 0.2449, \\
        \tfrac{\partial L}{\partial a} &= w_2\delta_2 \approx 0.4898, \qquad \delta_1 = \sigma'(z)\,\tfrac{\partial L}{\partial a} = a(1 - a) \cdot 0.4898 \approx 0.2350 \cdot 0.4898 \approx 0.1151, \\
        \tfrac{\partial L}{\partial w_1} &= \delta_1x \approx 0.1151, \qquad \tfrac{\partial L}{\partial b_1} = \delta_1 \approx 0.1151.
    \end{aligned}\]
    \textbf{Gradient check} with central differences, $h = 10^{-6}$: $\frac{L(w_1 + h) - L(w_1 - h)}{2h} \approx 0.11511$ and $\frac{L(w_2 + h) - L(w_2 - h)}{2h} \approx 0.15245$ \checkmark.
    \\[0.3em]
    Note the factor $\sigma'(z) \leq \frac{1}{4}$ in $\delta_1$: each sigmoid layer shrinks the signal (vanishing gradients).
}

\begin{frame}{Vanishing and exploding gradients}
    \sidebyside{0.5}{%
        \small
        $\boldsymbol{\delta}_1$ is $\boldsymbol{\delta}_L$ multiplied by the matrices
        \[\diag(\sigma'(\mathbf{z}_l))\,W_{l+1}^\top,\]
        once per layer.
        \begin{itemize}
            \item Sigmoid: $\sigma' \leq \frac{1}{4}$, so gradients of the first layers \textbf{vanish}.
            \item Large weights: gradients \textbf{explode}.
        \end{itemize}
        Remedies: ReLU ($\sigma' = 1$ on active units), careful initialization (Xavier/He, Lecture~9), residual connections $\mathbf{a}_l = \mathbf{a}_{l-1} + f_l(\mathbf{a}_{l-1})$ with Jacobian $I + Df_l$.
    }{%
        \centering\fig[1]{s07_vanishing}\\
        {\scriptsize random networks, $30$ layers of width $200$}
    }
    \vskip0.3em
    \small \textbf{Gradient checking:} compare with central differences; relative error $\frac{\norm{\mathbf{g} - \tilde{\mathbf{g}}}}{\norm{\mathbf{g}} + \norm{\tilde{\mathbf{g}}}} \approx 10^{-8}$ is fine, $\approx 1$ is a bug.
\end{frame}

% ------------------------------------------------------------------ Hessian
\topic{\secref{sec:hessian}}{Hessians and Taylor's theorem}

\begin{frame}{The Hessian matrix}
    \begin{definition}[][def:hessian]
        \small $\big(\nabla^2 f(\bx)\big)_{ij} = \frac{\partial^2 f}{\partial x_i\partial x_j}(\bx)$, the Jacobian of the gradient.
    \end{definition}
    \begin{theorem}[Schwarz][thm:schwarz]
        \small If $f$ is twice continuously differentiable, the mixed partials agree: the Hessian is symmetric.
    \end{theorem}
    \small
    \begin{example}[][ex:hessian]
        $f = x_1^2x_2 + x_1x_2^3$: \ $\nabla^2 f(\bx) = \begin{pmatrix} 2x_2 & 2x_1 + 3x_2^2 \\ 2x_1 + 3x_2^2 & 6x_1x_2 \end{pmatrix}$, \ $\nabla^2 f(1, 2) = \begin{pmatrix} 4 & 14 \\ 14 & 12 \end{pmatrix}$.
    \end{example}
    \begin{example}[Hessians in machine learning][ex:ml-hessians]
        $\frac{1}{2}\bx^\top A\bx - \mathbf{b}^\top\bx$: \ $\nabla^2 f = A$; \ least squares: $2\Phi^\top\Phi$; \ logistic regression: $\frac{1}{N}X^\top SX$ with $S = \diag(p_i(1 - p_i))$; \ log-sum-exp: $\diag(\mathbf{s}) - \mathbf{s}\mathbf{s}^\top$.
    \end{example}
\end{frame}

\begin{frame}{Taylor's theorem in several variables}
    \begin{theorem}[Second-order Taylor][thm:taylor2]
        \small For some $\tau \in (0, 1)$:
        \[f(\bx + \mathbf{h}) = f(\bx) + \nabla f(\bx)^\top\mathbf{h} + \tfrac{1}{2}\mathbf{h}^\top\nabla^2 f(\bx + \tau\mathbf{h})\,\mathbf{h},\]
        and hence $f(\bx + \mathbf{h}) = f(\bx) + \nabla f(\bx)^\top\mathbf{h} + \frac{1}{2}\mathbf{h}^\top\nabla^2 f(\bx)\,\mathbf{h} + o(\norm{\mathbf{h}}^2)$.
    \end{theorem}
    \small
    Proof idea: apply one-variable Taylor to $g(t) = f(\bx + t\mathbf{h})$; by the chain rule $g'(t) = \nabla f^\top\mathbf{h}$, $g''(t) = \mathbf{h}^\top\nabla^2 f\,\mathbf{h}$.
    \\[0.5em]
    \textbf{Preview of Lecture 8.} At a critical point ($\nabla f = \mathbf{0}$), $f(\bx + \mathbf{h}) \approx f(\bx) + \frac{1}{2}\mathbf{h}^\top\nabla^2 f(\bx)\mathbf{h}$, and the spectral theorem decides the shape: all eigenvalues $> 0$ a bowl, all $< 0$ a cap, mixed signs a saddle.
\end{frame}

\board{Example: a second-order approximation}{$f(x, y) = e^x\cos y$ at $(0, 0)$: gradient, Hessian, the quadratic approximation, and its value at $(0.1, 0.2)$ (exact: $1.0831\ldots$).}

\checkslide{$e^x\cos y$ near the origin}{
    \[\nabla f = \begin{pmatrix} e^x\cos y \\ -e^x\sin y \end{pmatrix}, \qquad \nabla^2 f = \begin{pmatrix} e^x\cos y & -e^x\sin y \\ -e^x\sin y & -e^x\cos y \end{pmatrix}\]
    At the origin: $f = 1$, $\nabla f = (1, 0)^\top$, $\nabla^2 f = \diag(1, -1)$:
    \[e^x\cos y \approx 1 + x + \tfrac{1}{2}(x^2 - y^2).\]
    At $(0.1, 0.2)$: linear $1 + 0.1 = 1.1$; quadratic $1.1 + \frac{1}{2}(0.01 - 0.04) = 1.085$; exact $e^{0.1}\cos 0.2 \approx 1.10517 \cdot 0.98007 \approx 1.0831$. \ Errors $0.017$ and $0.0019$. \\[0.4em]
    The Hessian $\diag(1, -1)$ has eigenvalues of both signs: the quadratic part is a saddle.
    \\[0.4em]
    \textbf{Another Hessian} (Example~\notesref{ex:hessian}): $f = x_1^2x_2 + x_1x_2^3$ at $(1, 2)$: $\partial_{11}f = 2x_2 = 4$, $\partial_{12}f = 2x_1 + 3x_2^2 = 14 = \partial_{21}f$, $\partial_{22}f = 6x_1x_2 = 12$: \ $\nabla^2 f(1, 2) = \begin{pmatrix} 4 & 14 \\ 14 & 12 \end{pmatrix}$, symmetric (Schwarz).
}

% ------------------------------------------------------------------ Gradient flow
\topic{\secref{sec:euler}}{Gradient flow and Euler's method}

\begin{frame}{Ordinary differential equations}
    \begin{definition}[][def:ode]
        \small A solution of $\bx' = F(\bx)$, $\bx(0) = \bx_0$, is a curve with $\bx'(t) = F(\bx(t))$ for all $t$.
    \end{definition}
    \begin{theorem}[Existence and uniqueness][thm:picard]
        \small If $F$ is Lipschitz, $\norm{F(\bx) - F(\by)} \leq L\norm{\bx - \by}$, there is exactly one solution for each $\bx_0$, defined for all $t$.
    \end{theorem}
    \small
    $x' = ax$: \ $x(t) = e^{at}x_0$. \ Linear systems $\bx' = A\bx$, $A = V\Lambda V^{-1}$: \ $\bx(t) = \sum_i c_ie^{\lambda_it}\bv_i$ (\S3.15.1).
    \begin{definition}[Gradient flow][def:gradient-flow]
        \small $\bx' = -\nabla f(\bx)$: always move in the direction of steepest descent.
    \end{definition}
    \begin{proposition}[][prop:gradient-flow]
        \small Along every solution, $\frac{d}{dt}f(\bx(t)) = -\norm{\nabla f(\bx(t))}^2 \leq 0$.
    \end{proposition}
\end{frame}

\board{Gradient flow of a quadratic}{Proof of the proposition (chain rule). \quad $f(\bx) = \frac{1}{2}\bx^\top A\bx - \mathbf{b}^\top\bx$, $A$ SPD, $\mathbf{e} = \bx - \bx^*$: show $\mathbf{e}' = -A\mathbf{e}$ and $\mathbf{e}(t) = \sum_i c_ie^{-\lambda_it}\bq_i$.}

\checkslide{gradient flow and Euler's method}{
    \small
    \textbf{Proof.} $\frac{d}{dt}f(\bx(t)) = \nabla f(\bx)^\top\bx'(t) = \nabla f^\top(-\nabla f) = -\norm{\nabla f}^2$.
    \\[0.3em]
    \textbf{Quadratic.} $\nabla f = A\bx - \mathbf{b} = A(\bx - \bx^*)$ since $A\bx^* = \mathbf{b}$; so $\mathbf{e}' = \bx' = -A\mathbf{e}$. Write $\mathbf{e} = \sum_ic_i(t)\bq_i$: each coordinate solves $c_i' = -\lambda_ic_i$, so $c_i(t) = c_i(0)e^{-\lambda_it}$.
    \\[0.2em]
    Example $A = \diag(1, 10)$, $\mathbf{e}(0) = (-2, 1)^\top$: \ $\mathbf{e}(t) = (-2e^{-t},\ e^{-10t})^\top$: the steep direction is gone at $t \approx 0.5$, the flat one needs $t \approx 5$.
    \\[0.4em]
    \textbf{Euler for $x' = -x$, $x(0) = 1$:} $x_{k+1} = x_k - hx_k = (1 - h)x_k$, so $x_k = (1 - h)^k$. At $t = 1$:
    \begin{center}
        \renewcommand{\arraystretch}{1.2}
        \begin{tabular}{@{}lccc@{}}
            \toprule
            & $h = 0.5$ & $h = 0.1$ & exact \\
            \midrule
            $x(1)$ & $0.5^2 = 0.25$ & $0.9^{10} \approx 0.349$ & $e^{-1} \approx 0.368$ \\
            error & $0.118$ & $0.019$ & \\
            \bottomrule
        \end{tabular}
    \end{center}
    Error $\propto h$. \ For $h > 2$: $|1 - h| > 1$, the iterates oscillate and blow up, although the true solution decays.
}

\begin{frame}{Gradient descent is Euler's method}
    \small
    Euler's method for $\bx' = F(\bx)$: follow the tangent for a time $h$, \ $\bx_{k+1} = \bx_k + hF(\bx_k)$. \\
    For the gradient flow: $\bx_{k+1} = \bx_k - h\nabla f(\bx_k)$, \ \textbf{gradient descent with $\eta = h$}.
    \\[0.3em]
    \centering\fig[0.78]{s07_euler}
    \\[0.1em]
    \raggedright
    Each step multiplies the error along $\bq_i$ by $1 - h\lambda_i$: stable iff $h < 2/\lambda_{\max}$; the slowest direction needs time $1/\lambda_{\min}$, so about $\kappa = \lambda_{\max}/\lambda_{\min}$ steps. The zig-zag is an instability of the discretization, not of the flow.
\end{frame}

\begin{frame}{ODEs in machine learning}
    \begin{mlbox}
        \small
        \begin{itemize}
            \item A residual network computes $\bx_{l+1} = \bx_l + f_l(\bx_l)$: Euler steps with $h = 1$. Neural ODEs replace the layers by an ODE solver, and get gradients from an adjoint ODE solved backwards in time (reverse mode in continuous time).
            \item Momentum is a discretization of the heavy ball $\bx'' + \gamma\bx' + \nabla f(\bx) = \mathbf{0}$.
            \item Adding noise gives stochastic differential equations: Langevin dynamics and diffusion models (Lecture~10, self-study \S4.15).
        \end{itemize}
    \end{mlbox}
    \vskip0.4em
    \textbf{Self-study} (\S3.13--3.14, \S3.15.1, \S3.16): vector fields, divergence and curl; multiple integrals, Fubini, change of variables, line integrals, the divergence theorem and integration by parts (score matching); linear ODE systems; how flows move probability densities (continuity equation, continuous normalizing flows).
\end{frame}

\begin{frame}{Self-study preview: vector fields and flows}
    \centering\notesfig[0.95]{L3-vector-fields}
    \\[0.3em]
    \raggedright\small
    The divergence $\Div F = \tr DF$ is the rate at which the flow of $F$ expands volumes. A density moved by the flow satisfies $\frac{\partial p_t}{\partial t} = -\Div(p_tF)$, and along a solution $\frac{d}{dt}\log p_t(\bx(t)) = -\Div F(\bx(t))$: continuous normalizing flows use this instead of a Jacobian determinant.
\end{frame}

% ------------------------------------------------------------------ Wrap-up
\begin{frame}{Key takeaways}
    \small
    \begin{itemize}
        \item Differentiable $=$ approximated by a linear map: $f(\bx + \mathbf{h}) = f(\bx) + Df(\bx)\mathbf{h} + o(\norm{\mathbf{h}})$. Rows of the Jacobian: outputs; columns: inputs.
        \item $Df(\bx)$ is a row vector (functional); $\nabla f(\bx)$ represents it. Steepest ascent and $\perp$ level sets: Cauchy--Schwarz and the chain rule.
        \item Chain rule: $D(f \circ g) = Df\,Dg$; \ $\nabla(f \circ g) = Dg^\top\nabla f$. Write the sizes.
        \item Matrix gradients: expand $f(X + H)$, read off $\ip{G}{H}$. Linear layer: $\nabla_W L = \boldsymbol{\delta}\bx^\top$, $\nabla_{\bx}L = W^\top\boldsymbol{\delta}$.
        \item Reverse mode (backprop) gives the whole gradient for a few evaluations of $f$; $\boldsymbol{\delta}_l = \sigma'(\mathbf{z}_l) \odot (W_{l+1}^\top\boldsymbol{\delta}_{l+1})$.
        \item Hessian: symmetric; $f(\bx + \mathbf{h}) \approx f + \nabla f^\top\mathbf{h} + \frac{1}{2}\mathbf{h}^\top\nabla^2 f\,\mathbf{h}$.
        \item Gradient descent is Euler's method for $\bx' = -\nabla f(\bx)$; stable only for $\eta < 2/\lambda_{\max}$.
    \end{itemize}
\end{frame}

\begin{frame}{Homework and next time}
    \small
    \begin{itemize}
        \item \textbf{Homework:} Lecture 3, Exercises 4, 8, 10, 11, 13 and 37; Programming Exercises P4 and P6.
        \item \textbf{Practice:} Practice 3, sections 1--8, and the gradient flow problems of section 15.
        \item \textbf{Self-study:} \S3.13--3.14 vector fields and multiple integrals; \S3.15.1 linear systems of ODEs; \S3.16 flows and densities.
    \end{itemize}
    \vskip0.6em
    \begin{block}{Next: Calculus II and optimization}
        Critical points and the second-order test; gradient descent and its convergence; momentum and Newton's method; convexity; the proximal gradient method and the lasso; Lagrange multipliers, KKT, duality.
    \end{block}
\end{frame}

\end{document}
